\scanpage{124}
\mspage{115}
% archive sheet number: 124

\[
\deleted{\Delta_{\mathrm{is}}=\unsure{(\Delta\setminus\partial\Psi)}\Delta\setminus\underbrace{\Delta\cap\mathrm{Omb}^{\circ}(\Psi_1)}_{\Delta_{\mathrm{int}}}}
\]

Posant
\[
\Delta_{\mathrm{int}}=\Delta\cap\mathrm{Omb}^{\circ}(\Psi_1)\subset\Delta\setminus\partial\Psi
\]
(\unsure{ens.} des pts de $\Delta$ \uline{intérieurs} \unsure{à} \deleted{\unsure{une figure}}\inserted{\unsure{l'un}} des intervalles $\unsure{S|\mathscr{S}}_{a_i,b_i}$) on aura
\[
\Delta=\Delta_{\mathrm{is}}\amalg\Delta_{\mathrm{int}}\amalg\partial\Psi
\]
où
\[
\Delta_{\mathrm{is}}\overset{\mathrm{def}}{=}(\Delta\setminus\partial\Psi)\setminus\Delta_{\mathrm{int}}
\]

\deleted{3}4°) Une partie
\[
\text{\textcircled{$\Delta_{\mathrm{lac}}$}}\subset\Delta_{\mathrm{int}}
\]
% slanted
\margin{\deleted{\unsure{déf.}} $\downarrow$}
et on \deleted{\unsure{pos}} aura donc
\[
\Delta_{\mathrm{int}}=\Delta_{\mathrm{lac}}\amalg\Delta_{\mathrm{red}}
\]
où
\[
\Delta_{\mathrm{red}}\overset{\mathrm{def}}{=}\Delta_{\mathrm{int}}\setminus\Delta_{\mathrm{lac}}\,.
\]

\deleted{\unsure{4°|3°})}\inserted{2°)} Une partie
\[
\text{\textcircled{$\unsure{\partial_0|\delta_0}$}}\subset\partial\Psi
\]
de $\partial\Psi$.

À ces données, on associe la préfigure $\Phi$ définie ainsi :
\[
\left\{\begin{aligned}
\deleted{\unsure{\delta}}\Phi_0&=\underbrace{\partial_0}_{\substack{\text{sommets}\\\text{bord propres}}}\cup\underbrace{\Delta_{\mathrm{red}}}_{\substack{\text{\unsure{sommet|sommets}}\\\text{redondants}}}\cup\Delta_{\mathrm{is}}\\
\Phi_1&=\{X\in\underbrace{\Phi_{\Delta}\cap\unsure{M|\mathcal{M}}_1}_{\substack{\text{intervalles interstitiels}\\\text{de }\Delta}}\mid X\in\mathrm{Omb}(\Psi)\}
\end{aligned}\right.
\]
